\chapter{Optimizing for Latency, Tokens, and Cost}
\label{ch:optimizing}

\epigraph{``Every second I'm in this room, I'm getting worse.''}{Frankie Dunn,
\emph{Million Dollar Baby} (2004)}

This is the chapter the rest of the book has been setting up. Everything from
here is a technique, a measurement, and an honest note about when it does not
help.

One ground rule before we start, and it is the whole discipline: \textbf{every
change gets attributed}. Not ``we did six things and it got faster''. Six
numbers, one per change, including the ones that turned out to be worth nothing.
Chapter~\ref{ch:measuring} gave you the instrument; this chapter uses it.

\section{The Meridian optimisation pass}

Same task throughout: score an account for credit risk, screen it for fraud, and
fetch the reference curve. Four servers, 12\,ms simulated per-call latency.

\bookfig[0.98]{optimization}{Four changes, measured cumulatively. Note that the
two panels improve at different steps, which is the point of measuring them
separately.}{optimization}

\begin{measurebox}{The full pass, attributed}
\begin{tabular}{@{}lrrrr@{}}
\toprule
\thead{Change} & \thead{Iterations} & \thead{Wall} & \thead{Tokens} & \thead{Cost} \\
\midrule
Baseline (38 tools, one call per turn) & 4 & 1{,}753\,ms & 22{,}846 & \$0.0690 \\
Trim the catalogue to 10 tools         & 4 & 1{,}751\,ms & 5{,}370  & \$0.0166 \\
Batch independent calls into one turn  & 2 & 1{,}067\,ms & 2{,}692  & \$0.0086 \\
Fan out that batch in parallel         & 2 & 1{,}014\,ms & 2{,}692  & \$0.0086 \\
\midrule
\textbf{Total} & \textbf{4 $\rightarrow$ 2} & \textbf{-42.2\,\%} & \textbf{-88.2\,\%} & \textbf{-87.6\,\%} \\
\bottomrule
\end{tabular}
\end{measurebox}

Read that table row by row, because each row teaches something different.

\textbf{Trimming the catalogue saved 76\,\% of tokens and zero milliseconds.}
The catalogue is not on the critical path for latency; it is on the critical
path for cost. Two budgets, two levers, and confusing them is how people spend a
sprint on the wrong one.

\textbf{Batching saved both, and by the most.} Four iterations became two, which
halved the model turns, which is where 96\,\% of the wall clock lives. This is
the round-trip lesson from Chapter~\ref{ch:primer}, cashed out.

\textbf{Parallel fan-out saved 53\,ms of a 1,067\,ms task.} Real, worth having,
and five percent. At 12\,ms per call there is not much to win; at 40\,ms there
is three times more.

\keyidea{Optimise iterations first, catalogue second, transport third. That is
the order of magnitude order, and it is stable across every system I have
measured.}

\section{Cut round trips}

The category that returns the most, always.

\subsection{Front-load MRTR inputs}

An elicitation costs about 375\,ms, of which 35 is transport and 340 is the
model turn it forces (Chapter~\ref{ch:mrtr}). So the cheapest elicitation is the
one that never happens.

\begin{py}
"exposureUsd": {
    "type": "number", "minimum": 0,
    "description": "Proposed exposure. Above 5,000,000 an "
                   "approver is required before scoring.",
},
\end{py}

That description is an optimisation. The model reads it, sees the threshold, and
can collect the approver from the user \emph{before} calling, in a turn it was
going to spend anyway. When it works, the round trip disappears entirely.

Better still, take the approver as an optional argument, so a model that already
knows it can supply it and skip the interrupt.

\subsection{Batch arguments}

\texttt{rank\_portfolio\_risk} exists for exactly one reason:

\begin{measurebox}{Portfolio review, 40 accounts}
\begin{tabular}{@{}lrr@{}}
\toprule
& \thead{Per-account tool} & \thead{Batch tool} \\
\midrule
Round trips           & 40 & 1 \\
Model turns           & 41 & 2 \\
Estimated wall clock  & \textasciitilde{}18\,s & \textasciitilde{}0.9\,s \\
\bottomrule
\end{tabular}

\vspace{6pt}
\footnotesize
Wall clock extrapolated from the measured 440\,ms mean model turn. The round-trip
counts are exact.
\end{measurebox}

Twenty times faster, from one extra tool. And the server-side description points
the model at it:

\begin{py}
"Rank a set of accounts by risk score in one call. Prefer this over "
"calling assess_account_risk repeatedly."
\end{py}

\subsection{Let the model batch, then fan out}

Two separate things that people conflate.

\emph{Batching} is the model emitting three tool calls in one turn instead of
three turns. That saves model turns, which is the expensive part.

\emph{Fan-out} is the host executing those three concurrently. That saves
transport time, which is the cheap part.

The measurement above shows the split precisely: batching took 1,751\,ms to
1,067\,ms, and fan-out took 1,067 to 1,014. Batching is thirteen times more
valuable, and it is the one that depends on your \emph{tool design}.

\subsection{Prefetch App templates}

Predeclared \texttt{ui://} templates can be fetched during idle time
(Chapter~\ref{ch:apps}). Meridian's is 2.6\,KB with a 24-hour public TTL, so
after the first fetch it is free forever. Cold, it is a resource read in the
critical path of the first render.

\section{Cache at every layer}

\subsection{Honour \texttt{ttlMs}}

The single cheapest win in the book, and the measurement is stark:

\begin{measurebox}{Per-task setup, 20 consecutive tasks, four servers}
\begin{tabular}{@{}lrr@{}}
\toprule
& \thead{Re-fetch every task} & \thead{Honour \texttt{ttlMs}} \\
\midrule
Requests issued (total)  & 164 & 8 \\
Setup time per task      & 274.1\,ms & 2.5\,ms \\
\bottomrule
\end{tabular}

\vspace{6pt}
\footnotesize
The naive host re-runs \texttt{server/discover} and \texttt{tools/list} against
all four servers before every task. The honouring host does it once.
\end{measurebox}

Twenty times fewer requests, and a hundred times less setup latency, from
respecting a field the server already sends you. If you implement one thing from
this chapter, implement this.

And the resource-read side, from Chapter~\ref{ch:resources}: 577 of 600 reads
avoided, 96.2\,\% hit rate, 6.82$\times$ faster.

\subsection{Use \texttt{cacheScope: public} where it is true}

A \texttt{public} result can be served by a shared gateway to every user. For
Meridian's market-data server, which returns identical data to everybody, that
turns N clients' worth of load into one fetch.

For anything personalised, it is a data breach. Chapter~\ref{ch:resources} has
the reasoning; the rule is that \texttt{public} is a claim about the
\emph{content}, not about the endpoint's authentication.

\subsection{Stable ordering, for the provider's prompt cache}

Deterministic \texttt{tools/list} ordering is a specification SHOULD, and the
reason is prompt caching. Providers cache on a stable prefix. Reorder your
catalogue and every request misses.

\begin{py}
def test_tools_list_order_is_stable(self):
    """Unstable ordering silently destroys the provider's prompt cache."""
\end{py}

The word ``silently'' is the problem. Nothing errors. Your server looks healthy.
You just pay full prompt-processing price on every turn forever.

The same discipline applies to your host's merged catalogue
(\texttt{sorted(self.bindings)}) and to your system prompt, which should not
contain a timestamp.

\subsection{Invalidate on notification}

TTL and \texttt{listChanged} together let you run long TTLs safely. Without the
notification you must choose between staleness and refetching; with it you can
have a five-minute TTL and still see changes in milliseconds.

\section{Stream everything}

Perceived latency is the latency. The Meridian baseline is 1,366\,ms and 96\,\%
of it is model inference you cannot speed up, so the remaining lever is making
the wait legible.

\begin{itemize}[leftmargin=1.6em, itemsep=3pt]
\item \textbf{Token streaming.} Time to first token is 340\,ms of a 440\,ms
turn. Stream, and the user waits 340\,ms instead of 440.
\item \textbf{Progress notifications.} Free UX from any server that emits them.
Forward them.
\item \textbf{Partial results.} Chapter~\ref{ch:agent-loop} argued for this;
the measurement here is why. Iteration 0 finishes at 447\,ms of a 1,366\,ms
task, so showing its result buys the user two thirds of the wait back.
\item \textbf{\texttt{X-Accel-Buffering: no}} on every SSE response, or a proxy
will buffer your carefully streamed events into one lump at the end.
\end{itemize}

\section{Transport tuning}

The smallest category, included because it is cheap and because people ask.

\subsection{Reuse connections}

\begin{py}
def test_sequential_calls_reuse_one_connection(self):
    transport = StreamableHttpClient(self.server.url)
    client = Client(transport, server_label="wire")
    for _ in range(5):
        client.call_tool("echo", {"text": "x"})
    self.assertEqual(transport.new_connections, 1)
    self.assertGreaterEqual(transport.reused_connections, 4)
    transport.close()
\end{py}

On loopback this saves almost nothing. Over TLS to another region it saves a TCP
handshake plus a TLS handshake per call, which is two extra round trips, which at
68\,ms each is 136\,ms you were about to pay repeatedly.

\subsection{One connection is not enough, and this one bit me}

Reuse the connection, and then notice what you have built. HTTP/1.1 carries one
request at a time per connection: the next request cannot start until the
previous response has been read. That is head-of-line blocking, and it is why a
browser opens six connections per origin.

Now put that under the parallel fan-out this chapter recommends. A host that
issues three tool calls concurrently, all to the same server, through one warm
keep-alive connection, gets three requests queued behind each other. The fan-out
code runs, the threads start, and the wall clock is the sum. The optimisation is
present, measurable in the code, and worth zero.

Meridian had exactly this shape, and the book's 2.8$\times$ fan-out figure hid it
in plain sight, because that measurement fans out across \emph{different}
servers, each with its own client and its own connection. Every number in this
book is honest and the reader fanning out to one server would have got nothing.

The fix is a pool:

\begin{py}
def _acquire(self):
    """Take an idle connection, or make one.

    A pool rather than a single connection, because HTTP/1.1 has no
    multiplexing: one connection carries one request at a time, so three
    parallel tool calls sharing one connection run in sequence and the
    fan-out buys nothing. The pool is what makes the host's parallel path
    parallel when every call goes to the same server.
    """
\end{py}

Bounded, because unbounded is its own failure: a host that opens a connection per
concurrent call will exhaust file descriptors on the client and connection slots
on the server, and does so first during the traffic spike you built the fan-out
for. Six is the number browsers settled on and it is a reasonable default.

\begin{notebox}{HTTP/2 makes this disappear}
HTTP/2 multiplexes many streams over one connection, so the pool becomes
unnecessary and the head-of-line blocking moves down to the TCP layer, where a
lost packet stalls every stream sharing the connection. HTTP/3 fixes that in
turn by running over QUIC.

Meridian is HTTP/1.1 because it is built on the standard library and the point
was to show the protocol rather than the plumbing. If your deployment already
speaks HTTP/2 end to end, you get the concurrency for free and can set the pool
to one.
\end{notebox}

\subsection{Warm pools for stdio}

\begin{measurebox}{stdio cold start}
\begin{tabular}{@{}lr@{}}
\toprule
Spawn to first usable response & 44.6\,ms \\
Warm call                      & 0.060\,ms \\
Calls to amortise one spawn    & \textasciitilde{}750 \\
\bottomrule
\end{tabular}
\end{measurebox}

Spawn-per-call is indefensible for anything but a one-shot CLI. Keep a pool,
size it to your concurrency, and recycle processes on a schedule.

\subsection{Let the gateway route on headers}

\texttt{Mcp-Method} and \texttt{Mcp-Name} exist so intermediaries need not parse
bodies. Use them: send \texttt{tools/list} to a cache tier and
\texttt{tools/call} to a compute tier, rate-limit per method, apply strict WAF
policy to the two tools that matter and relaxed policy to the rest.

And validate them, always, because two sources of truth that can disagree is a
request-smuggling primitive (Chapter~\ref{ch:transport}).

\section{Spend tokens deliberately}

\subsection{Right-size the catalogue}

The biggest single token lever, measured twice now:

\begin{measurebox}{Catalogue cost}
\begin{tabular}{@{}lrr@{}}
\toprule
& \thead{REST mirror} & \thead{Task-shaped} \\
\midrule
Tools                     & 38 & 10 \\
Tokens per turn           & 5{,}617 & 1{,}246 \\
Tokens per task (4 turns) & 22{,}846 & 5{,}370 \\
Cost per task             & \$0.0690 & \$0.0166 \\
\bottomrule
\end{tabular}
\end{measurebox}

Four techniques, in order of how much they usually return:

\begin{enumerate}[leftmargin=1.6em, itemsep=3pt]
\item \textbf{Delete unused tools.} Check your telemetry. Anything with no calls
in 90 days is pure cost.
\item \textbf{Merge near-duplicates.} Three tools that differ by one flag should
be one tool with a parameter.
\item \textbf{Compress descriptions.} Lead with the outcome, drop the transport
trivia. Chapter~\ref{ch:tools} has the before and after.
\item \textbf{Filter per task.} Send the tools plausibly relevant to this
request. Costs a retrieval step, and it is the only technique that scales past a
few dozen tools.
\end{enumerate}

\subsection{Structured outputs eliminate parse-retry loops}

Chapter~\ref{ch:tools} made the correctness argument. The performance argument
is the retry: a misparse costs a full model turn, roughly 440\,ms and 1,300
input tokens, and client-side validation against an \texttt{outputSchema} catches
the malformed response at the boundary instead.

An \texttt{enum} is worth extra attention. Telling the model that \texttt{band}
is one of four values stops it inventing a fifth and stops it writing defensive
handling for cases that cannot occur.

\subsection{Truncate tool results before they enter context}

\begin{py}
return text if len(text) <= limit else text[:limit] + "...[truncated]"
\end{py}

Blunt, and better than unbounded. Chapter~\ref{ch:agent-loop} has the more
sophisticated options; the important thing is that a bound exists, because a
single 40\,KB tool result poisons every subsequent turn of the task.

\subsection{Tune task polling}

For the Tasks extension, cadence is a direct token and request cost. Honour
\texttt{pollIntervalMs}, back off for long jobs, and add jitter so a thousand
clients that started together do not stay together.

\section{Route models}

Not every turn needs your best model.

\begin{table}[htbp]
\centering
\small
\begin{tabularx}{\textwidth}{@{}lXl@{}}
\toprule
\thead{Turn} & \thead{Cognitive load} & \thead{Model} \\
\midrule
Initial planning            & High. Real reasoning & Large \\
Tool selection from 3 options & Low. Classification & Small \\
Reading a structured result & Low. Extraction & Small \\
Deciding whether to stop    & Medium. Judgement & Large \\
Summarising for context     & Low. Compression & Small \\
Final synthesis             & High. Composition & Large \\
\bottomrule
\end{tabularx}
\caption{In a long task the middle rows are the majority of turns, which is what
makes routing worth the complexity.}
\label{tab:routing}
\end{table}

Two warnings, both learned expensively.

\textbf{Measure the frontier, do not reason about it.} Run your eval suite at
each mix. It is common to move most turns to a small model and lose nothing
measurable. It is also common to lose 15 points of task success, and which one
you are in is not predictable from first principles.

\textbf{A cheap wrong answer is not cheap.} A small model that picks the wrong
tool costs a full recovery cycle: another model turn, another tool call, another
model turn. That is more expensive than the large model would have been.

\section{The order to do things in}

\begin{table}[htbp]
\centering
\small
\begin{tabularx}{\textwidth}{@{}lXll@{}}
\toprule
\thead{\#} & \thead{Change} & \thead{Typical return} & \thead{Effort} \\
\midrule
1 & Honour \texttt{ttlMs}              & 20$\times$ fewer setup requests & one afternoon \\
2 & Trim the catalogue                 & 76\,\% of tokens & one conversation \\
3 & Add batch variants of looped tools & 20$\times$ on portfolio tasks & one tool \\
4 & Batch and fan out                  & 40\,\% of wall clock & host change \\
5 & Deterministic ordering             & prompt-cache hits & one \texttt{sorted()} \\
6 & Stream tokens and progress         & perceived latency & UI work \\
7 & Connection reuse, warm pools       & 2 RTT per call & config \\
8 & Route models                       & 30 to 60\,\% of cost & eval work \\
9 & Rewrite the server in a fast language & 3\,ms & three weeks \\
\bottomrule
\end{tabularx}
\caption{Roughly descending value per unit effort. Row nine is from the preface,
and it is there to be looked at again.}
\label{tab:optimisation-order}
\end{table}

\section{What did not help}

The section every performance chapter should have and almost none do.

\textbf{Trimming the catalogue did not reduce latency.} Zero milliseconds, twice
measured. Tokens and time are different budgets.

\textbf{Caching did not speed up the loop itself.} Meridian fetches its catalogue
before the loop starts, so the loop never touches the cache. Caching helped the
setup path enormously and the loop not at all. If you had measured only the loop,
you would have concluded caching was useless.

\textbf{Parallel fan-out returned 5\,\% at 12\,ms per call.} Real, worth keeping,
and not the story. At 40\,ms it is worth three times more, which is a reminder
that the value of an optimisation depends on the environment you measured it in.

\textbf{The \texttt{\_meta} envelope's 160\,\% byte overhead is irrelevant.} It
looks alarming and it is 197 bytes on a request whose unit of work is hundreds of
milliseconds of inference. Do not optimise it. I have watched somebody try.

\keyidea{Publish the changes that did not work. They are how you stop the next
person spending a sprint on the thing you already measured.}

\section{Meridian, before and after}

\begin{measurebox}{The complete pass}
\begin{tabular}{@{}lrrr@{}}
\toprule
& \thead{Before} & \thead{After} & \thead{Change} \\
\midrule
Loop iterations       & 4 & 2 & -50\,\% \\
Wall clock            & 1{,}753\,ms & 1{,}014\,ms & -42.2\,\% \\
Transport time        & 104.5\,ms & 36.8\,ms & -64.8\,\% \\
Tokens per task       & 22{,}846 & 2{,}692 & -88.2\,\% \\
Cost per task         & \$0.0690 & \$0.0086 & -87.6\,\% \\
Setup requests (20 tasks) & 164 & 8 & -95.1\,\% \\
\midrule
Cost per 1,000 tasks  & \$69.01 & \$8.56 & \textbf{-\$60.45} \\
\bottomrule
\end{tabular}

\vspace{6pt}
\footnotesize
Regenerate with \texttt{make bench}. Every row is attributed to a specific
change in the table at the top of this chapter.
\end{measurebox}

No new hardware. No faster language. Four changes, each measured, each
attributable, and the largest of them was deleting tools.

\section{What to remember}

\begin{itemize}[leftmargin=1.6em, itemsep=3pt]
\item Attribute every change. Six numbers, not one story.
\item Iterations first, catalogue second, transport third.
\item Batching saves model turns; fan-out saves transport. Batching is worth
roughly thirteen times more.
\item Honouring \texttt{ttlMs} cut setup requests by 95\,\%. Do this first.
\item Deterministic ordering protects the provider's prompt cache, and losing it
is silent.
\item Front-load MRTR inputs. The cheapest elicitation is the one that never
happens.
\item Structured outputs remove parse-retry loops; enums remove invented values.
\item Route glue turns to small models, and verify with evals rather than
intuition.
\item A cheap wrong answer costs a full recovery cycle.
\item Publish what did not work.
\end{itemize}

Next: running it, once it is fast.
